\textbf{Desarrollo de las magnitudes intermedias.} De $D=W\sin\delta$ y $L=W\cos\delta$, dividir da $\tan\delta=C_D/C_L$. La sustentación exacta usa el coseno; el descenso usa el seno. La aproximación tendida sustituye ambos factores por sus límites de ángulo pequeño.\begin{enumerate}

\item Calcular la polar al CL impuesto.

\[\fitmath{C_D=0.014+0.032\cdot {0.78}^{2}\approx 0.0334688\,}\]

\item Obtener ángulo de descenso positivo.

\[\fitmath{\delta=\arctan\left(\frac{C_D}{0.78}\right)}\]

\[\fitmath{\delta=\arctan\left(\frac{(0.0334688)}{0.78}\right)\approx 0.04288241\,\mathrm{rad}}\]

\item Conservar coseno en el equilibrio normal.

\[\fitmath{\mathrm{Vexact}=\sqrt{\frac{2\cdot 455\cdot 9.80665\cdot \cos\left(\delta\right)}{1.05\cdot 16.2\cdot 0.78}}}\]

\[\fitmath{\mathrm{Vexact}=\sqrt{\frac{2\cdot 455\cdot 9.80665\cdot \cos\left((0.04288241309)\right)}{1.05\cdot 16.2\cdot 0.78}}\approx 25.9228\,\mathrm{m/s}}\]

\item Proyectar la velocidad sobre la vertical.

\[\fitmath{\mathrm{sexact}=\mathrm{Vexact}\cdot \sin\left(\delta\right)}\]

\[\fitmath{\mathrm{sexact}=(25.92279697)\cdot \sin\left((0.04288241309)\right)\approx 1.111291\,\mathrm{m/s}}\]

\item Aplicar solo en esta ruta el equilibrio aproximado L=W.

\[\fitmath{\mathrm{Vapprox}=\sqrt{\frac{2\cdot 455\cdot 9.80665}{1.05\cdot 16.2\cdot 0.78}}\approx 25.93472\,\mathrm{m/s}}\]

\item Sustituir seno por tangente en la aproximación.

\[\fitmath{\mathrm{sapprox}=\frac{\mathrm{Vapprox}\cdot C_D}{0.78}}\]

\[\fitmath{\mathrm{sapprox}=\frac{(25.93472073)\cdot (0.0334688)}{0.78}\approx 1.112826\,\mathrm{m/s}}\]

\item Cuantificar el error en descenso.

\[\fitmath{\mathrm{relative}=\frac{\mathrm{sapprox}}{\mathrm{sexact}}-1}\]

\[\fitmath{\mathrm{relative}=\frac{(1.112825617)}{(1.111291422)}-1\approx 0.001380551\,}\]

\end{enumerate}
